
\subsubsection{4.1 Research Questions}

\begin{itemize}
\item \textbf{RQ1:} Does task-relevant structure exist in pretrained hidden states?
\item \textbf{RQ2:} Can self-supervised embeddings capture task-discriminative patterns?
\item \textbf{RQ3:} Can low-rank modulation efficiently condition SSM dynamics?
\item \textbf{RQ4:} Does TC-SSM preserve adaptation capability?
\end{itemize}

\subsubsection{4.2 Datasets}

SuperGLUE validation splits:

\begin{table}[htbp]
\centering
\resizebox{\columnwidth}{!}{%
\begin{tabular}{lll}
\toprule
Task & Samples & Description \\
\midrule
BoolQ & 3,270 & Boolean QA \\
CB & 57 & CommitmentBank \\
COPA & 100 & Causal reasoning \\
RTE & 277 & Textual entailment \\
WiC & 638 & Word-in-context \\
WSC & 104 & Winograd schema \\
\bottomrule
\end{tabular}
}
\end{table}

For clustering experiments (h-e1), 660 samples were used across 5 tasks. For embedding learning (h-m1), 7,204 samples across 6 tasks (5,042 train, 1,082 test).

\subsubsection{4.3 Baselines}

\begin{itemize}
\item \textbf{Transformer baseline:} Fine-tuned classifier without architecture conversion
\item \textbf{Standard distillation:} Knowledge distillation without task conditioning
\end{itemize}

\subsubsection{4.4 Implementation}

\begin{itemize}
\item \textbf{Task discovery:} K-means (K=8) on BERT-base [CLS] embeddings, seed=42
\item \textbf{Embedding learning:} InfoNCE with $\tau =0.1$, AdamW (lr=1e-4), 10 epochs
\item \textbf{Modulation experiment:} d\_model=1024, d\_state=16, batch\_size=8, seq\_len=128
\item \textbf{Few-shot evaluation:} 8-shot and 16-shot protocols, 3 seeds, maximum 100 gradient steps
\end{itemize}

\subsubsection{4.5 Evaluation Metrics}

\begin{itemize}
\item Cluster purity, NMI, ARI for task discovery
\item Linear probe accuracy for embedding quality
\item FLOPs overhead and state variance ANOVA for modulation efficiency
\item Few-shot accuracy gap and steps to 95\% ceiling for adaptation preservation
\end{itemize}

